Un skieur de masse $M=\qty{80}{\kg}$ (avec son équipement) descend une pente
inclinée de l'angle $\alpha=\ang{10}$ par rapport à l'horizontale.

Il progresse dans la neige poudreuse et la force de frottement qui s'exerce sur
lui est constante, parallèle à son mouvement, en sens inverse de celui-ci et de
valeur égale à $f=\qty{60}{\N}$.

Le skieur étant initialement immobile, quelle est la valeur de sa vitesse après
un parcours de $\qty{100}{\m}$ sur cette pente $(g=\qty{9.8}{\m\per\s})$~?
